% ============================================================
%  Exposición -- Herramientas Automatizadas para la IR
%  Formato: IEEE onecolumn | Todo en negro
%  Compilación: pdflatex → biber → pdflatex → pdflatex
% ============================================================
\documentclass[journal,onecolumn]{IEEEtran}

\usepackage[T1]{fontenc}
\usepackage[utf8]{inputenc}
\usepackage[spanish,es-tabla]{babel}

\usepackage{booktabs}
\usepackage{tabularx}
\usepackage{array}
\usepackage{multirow}
\usepackage{longtable}

\usepackage{graphicx}
\usepackage{float}

\usepackage{xcolor}
\usepackage[hidelinks,colorlinks=false]{hyperref}

\usepackage{geometry}
\geometry{a4paper, top=2.5cm, bottom=2.5cm, left=2.5cm, right=2.5cm}

\usepackage{fancyhdr}
\pagestyle{fancy}
\fancyhf{}
\renewcommand{\headrulewidth}{0.4pt}
\fancyhead[L]{\small Ingeniería de Requisitos -- Herramientas Automatizadas}
\fancyhead[R]{\small MundiPets}
\fancyfoot[C]{\thepage}

\usepackage{parskip}
\setlength{\parskip}{4pt}
\usepackage{enumitem}
\usepackage{caption}
\usepackage{titlesec}
\usepackage[autostyle=true,spanish=mexican]{csquotes}
\captionsetup{font=small, labelfont=bf}
\renewcommand{\thetable}{\arabic{table}}

\usepackage{mdframed}
\mdfdefinestyle{guia}{
  linecolor=black,
  linewidth=0.6pt,
  backgroundcolor=gray!10,
  innertopmargin=6pt,
  innerbottommargin=6pt,
  innerleftmargin=8pt,
  innerrightmargin=8pt,
  skipabove=4pt,
  skipbelow=4pt
}
\newcommand{\guia}[2]{%
  \begin{mdframed}[style=guia]
    \noindent\textbf{Responsable: #1}\\[2pt]
    \textit{#2}
  \end{mdframed}
}

\usepackage[
  backend=biber,
  style=ieee,
  citestyle=numeric-comp,
  isbn=true,
  doi=true
]{biblatex}
\addbibresource{referencias_herramientas.bib}

\renewcommand{\thesection}{\arabic{section}}
\renewcommand{\thesubsection}{\thesection.\arabic{subsection}}
\renewcommand{\thesubsubsection}{\thesubsection.\arabic{subsubsection}}

\titleformat{\section}{\normalfont\Large\bfseries}{\thesection}{1em}{}
\titleformat{\subsection}{\normalfont\large\bfseries}{\thesubsection}{1em}{}
\titleformat{\subsubsection}{\normalfont\normalsize\bfseries}{\thesubsubsection}{1em}{}


% ============================================================
\begin{document}
	\begin{titlepage}
		\centering
		
		{\large\bfseries UNIVERSIDAD TÉCNICA ESTATAL DE QUEVEDO\\[12pt]}
		
		\includegraphics[width=0.3\textwidth]{uteqlogo.png}\\[12pt]
		
		{\large\bfseries FACULTAD DE CIENCIAS DE LA COMPUTACIÓN Y DISEÑO DIGITAL\\[4pt]
			CARRERA DE INGENIERÍA EN SOFTWARE}
		
		\vspace{1cm}
		
		{\normalsize\bfseries TEMA:}\\[4pt]
		{\normalsize EXPOSICIÓN - HERRAMIENTAS AUTOMATIZADAS PARA LA INGENIERÍA DE REQUISITOS}
		
		\vspace{1cm}
		
		{\normalsize \bfseries MATERIA:}\\[4pt]
		{\normalsize INGENIERÍA DE REQUERIMIENTOS}
		
		\vspace{1cm}
		
		{\normalsize\bfseries DOCENTE:}\\[4pt]
		{\normalsize ING. GUERRERO ULLOA GLEISTON CICERÓN}
		
		\vspace{1cm}
		
		{\normalsize\bfseries INTEGRANTES:}\\[6pt]
		
		\textit{\textbf{FUERTES ARRAES EDSON DANIEL}},~0929115087,~efuertesa2@uteq.edu.ec\\
		
		\vspace{0.5cm}
		
		\textit{CEDEÑO AVILA WINSTON DAMIAN},~0942833492,~wcedenoa2@uteq.edu.ec\\
		
		\vspace{0.5cm}
		
		\textit{CEDEÑO CORONADO WILSON LIZANDRO},~1206867209,~wcedenoc2@uteq.edu.ec
		
		\vspace{1cm}
		
		{\normalsize\bfseries CURSO:}\\[4pt]
		{\normalsize CUARTO SEMESTRE PARALELO ``B''}
		
		\vspace{1cm}
		
		{\normalsize\bfseries QUEVEDO -- LOS RÍOS -- ECUADOR}\\[4pt]
		{\normalsize\bfseries 2026 -- 2027 PPA}\\[4pt]
		{\normalsize Semana 6 -- \today}
		
	\end{titlepage}

\tableofcontents
\newpage

% ============================================================
\section{Introducción}
Las herramientas automatizadas de soporte a la Ingeniería de
Requisitos surgieron como respuesta a las limitaciones que presentan
los métodos manuales, como hojas de cálculo y documentos de texto, cuando
los proyectos crecen en número de requisitos, partes interesadas y
ciclos de cambio. Jallow et al.\ demuestran empíricamente que la
gestión manual de requisitos es intensiva en papel, carece de un
repositorio centralizado y hace prácticamente imposible mantener la
trazabilidad y el control de cambios a lo largo del ciclo de vida de
un proyecto \cite{jallow2014}. A medida que los sistemas de software
aumentan en complejidad, esta situación se agrava: los equipos trabajan
con versiones desactualizadas de los requisitos, los cambios se comunican por correo electrónico o verbalmente, y el análisis de impacto se realiza sin un proceso definido, con base en la memoria y la experiencia individual \cite{jallow2014}.

Frente a este panorama, las herramientas automatizadas de gestión de
requisitos aportan cuatro capacidades esenciales que los métodos
manuales no pueden escalar: trazabilidad bidireccional entre necesidades
de negocio, diseño, código y pruebas; control de versiones de cada
requisito a lo largo del tiempo; colaboración simultánea de múltiples
\textit{stakeholders} distribuidos geográficamente; y generación
automática de informes de impacto ante cambios. Siddique señala que la
tendencia actual en IR apunta precisamente hacia la automatización e
integración de estas actividades con el resto del ciclo de vida del
desarrollo de software, incorporando tecnologías como inteligencia
artificial, aprendizaje automático y procesamiento de lenguaje natural
\cite{siddique2023}. En esa misma línea, Groen et al.\ evidencian que
el análisis manual de grandes volúmenes de retroalimentación de usuarios
no escala bien y que la automatización reduce el tiempo de procesamiento
en un orden de magnitud \cite{groen2018}.

El objetivo de este documento es investigar, evaluar y comparar las
herramientas representativas del mercado de soporte automatizado a la
IR, analizando sus funcionalidades principales, ventajas,
limitaciones y aplicabilidad al sistema del Proyecto Fin de Curso (PFC).
Las herramientas fueron seleccionadas cubriendo espectros distintos:
desde soluciones enterprise de alto costo orientadas a industrias
reguladas, hasta alternativas ligeras y gratuitas adecuadas para equipos
académicos, incluyendo además la categoría emergente de la IA
generativa aplicada a la elicitación.

\subsection{Criterios de Selección de las Herramientas}

Para este estudio, ZenTao fue asignada por el docente como herramienta a investigar y exponer.
La selección responde a criterios que Wiegers y Beatty identifican como determinantes al evaluar
una herramienta de gestión de requisitos: costo total de propiedad, curva de aprendizaje,
capacidades de trazabilidad y facilidad de integración con el proceso de desarrollo del equipo
\cite{wiegers2013}. ZenTao cubre estos criterios en el contexto académico al ofrecer una edición
\textit{open source} sin costo de licencia, con módulos de \textit{backlog}, pruebas y defectos
integrados en una misma plataforma, y con una curva de aprendizaje manejable para equipos pequeños.
\newpage

% ────────────────────────────────────────────────────────────
\section{Análisis de Herramienta ZenTao}
% ────────────────────────────────────────────────────────────

\subsection{Descripción General}

ZenTao es una herramienta de gestión de proyectos de software que integra funcionalidades para
la administración de requisitos, planificación, pruebas y seguimiento de defectos dentro de una
misma plataforma \cite{zentao2025}. Está disponible en modalidad \textit{open source}, lo que la
convierte en una alternativa accesible para pequeñas y medianas organizaciones, así como para
entornos académicos. Desde la perspectiva de la Ingeniería de Requisitos, ZenTao se alinea con
conceptos fundamentales descritos en el SWEBOK V4 y en la norma ISO/IEC/IEEE 29148:2018,
particularmente en lo que respecta a la trazabilidad, la gestión del cambio y la colaboración
entre \textit{stakeholders} \cite{iso29148}.

\subsection{Funcionalidades Principales}

\begin{itemize}
  \item \textbf{Gestión de requisitos e historias de usuario:} permite registrar, organizar y
  priorizar requisitos como historias de usuario dentro de un \textit{backlog}, facilitando su
  seguimiento a lo largo del ciclo de vida del proyecto \cite{zentao2025}.

  \item \textbf{Trazabilidad:} relaciona requisitos con tareas de desarrollo, casos de prueba y
  defectos reportados, lo que permite rastrear cada requisito desde su origen hasta su verificación,
  en concordancia con lo establecido en la norma ISO/IEC/IEEE 29148 \cite{iso29148}.

  \item \textbf{Gestión del cambio:} mantiene un historial de modificaciones sobre cada tarea o
  requisito, lo que facilita el control de versiones y el análisis de impacto ante cambios durante
  el desarrollo \cite{zentao2025}.

  \item \textbf{Gestión de pruebas:} integra casos de prueba directamente vinculados a los
  requisitos, lo que apoya las actividades de validación descritas en el SWEBOK V4 \cite{swebok2024}.

  \item \textbf{Colaboración entre equipos:} centraliza la información del proyecto, permitiendo
  que analistas, desarrolladores y responsables de pruebas accedan a los mismos datos, reduciendo
  inconsistencias en la documentación y mejorando la comunicación \cite{zentao2025}.
\end{itemize}

\subsection{Ventajas}

\begin{itemize}
  \item \textbf{Centralización de la información:} al reunir requisitos, tareas, pruebas y
  defectos en una sola plataforma, se reduce la posibilidad de que información relevante quede
  desactualizada o dispersa entre los miembros del equipo \cite{zentao2025}.

  \item \textbf{Accesibilidad económica:} su modelo \textit{open source} implica menores costos
  de implementación en comparación con herramientas empresariales especializadas, lo que lo hace
  viable para proyectos académicos y organizaciones de pequeña y mediana escala \cite{zentao2025}.

  \item \textbf{Alineación con estándares de IR:} las funcionalidades de trazabilidad y gestión
  del cambio se corresponden directamente con las recomendaciones de la norma ISO/IEC/IEEE 29148
  y los principios descritos en el SWEBOK V4 \cite{iso29148, swebok2024}.
\end{itemize}

\subsection{Limitaciones}

\begin{itemize}
  \item \textbf{Curva de aprendizaje inicial:} Wiegers y Beatty advierten que introducir una
  herramienta de gestión de requisitos en un proyecto ya en marcha implica riesgos que suelen
  subestimarse, entre ellos la resistencia del equipo y los problemas que solo afloran en uso
  real \cite{wiegers2013}. En ZenTao, la coexistencia de módulos de \textit{backlog},
  \textit{sprints}, pruebas y defectos dentro de una misma plataforma puede resultar confusa
  para usuarios sin experiencia previa en herramientas ALM \cite{zentao2025}.

  \item \textbf{Menor especialización en IR formal:} Hussain et al.\ constatan que las
  herramientas ALM de propósito general presentan brechas frente a herramientas dedicadas
  exclusivamente a IR cuando se trata de gestionar trazabilidad a nivel de sistema en proyectos
  de gran escala o en industrias reguladas, donde se exigen matrices de trazabilidad complejas
  y auditorías de conformidad \cite{hussain2022}. ZenTao no ofrece, de forma nativa, los niveles
  de trazabilidad de sistema completos que exigen normas como DO-178C o IEC~62304.

  \item \textbf{Funcionalidades avanzadas restringidas a ediciones de pago:} características
  como la integración con sistemas de CI/CD, el tablero avanzado de métricas y el soporte
  oficial en múltiples idiomas están disponibles únicamente en las ediciones Enterprise o
  Professional, lo que puede representar una barrera para equipos académicos o de pequeña
  escala que operan con la versión \textit{open source} \cite{zentao2025}.
\end{itemize}


% ============================================================
\section{Análisis de Herramienta ReqView}
% ============================================================

\subsection{Descripción General}

ReqView es una herramienta de escritorio para la gestión de requisitos desarrollada por la
empresa Eccam s.r.o., fundada en 2004 en Praga, República Checa \cite{reqview2026}. Está
orientada principalmente a equipos pequeños y medianos que necesitan una solución sencilla
para organizar, documentar y dar seguimiento a los requisitos de un proyecto
\cite{khoroshilov2019}.

A diferencia de otras herramientas empresariales más complejas, ReqView se caracteriza por
ser una opción ligera y de menor costo, lo que la convierte en una alternativa atractiva
para pequeñas empresas, equipos de desarrollo y proyectos académicos. Estudios previos ya
destacaban que la herramienta ofrecía las funciones básicas necesarias para la gestión de
requisitos \cite{khoroshilov2019}.

Una de sus características más importantes es que almacena la información en archivos
locales con formato JSON. Esto permite a los usuarios mantener el control de sus datos,
trabajar sin conexión a Internet y utilizar herramientas de control de versiones como Git
para gestionar cambios y colaborar con otros miembros del equipo \cite{khoroshilov2019}.
Además, ReqView es compatible con los sistemas operativos Windows, Linux y macOS, y ofrece
una licencia comercial con un período de prueba gratuito de 14 días, permitiendo a los
usuarios evaluar sus funcionalidades antes de adquirir el producto.

\subsection{Funcionalidades Principales}

\begin{itemize}
	\item \textbf{Catálogo de requisitos:} permite organizar los requisitos mediante
	identificadores y atributos básicos, así como gestionar requisitos de
	\textit{stakeholders}, requisitos del sistema, riesgos y actividades de validación
	y verificación \cite{khoroshilov2019}.

	\item \textbf{Gestión de relaciones y trazabilidad:} permite establecer relaciones
	entre requisitos para dar seguimiento a su evolución y mantener la trazabilidad
	dentro del proyecto, aunque con opciones más limitadas que las de herramientas
	empresariales más avanzadas \cite{khoroshilov2019}.

	\item \textbf{Importación y exportación de información:} es compatible con formatos
	como Microsoft Word, Excel, PDF, HTML y ReqIF, facilitando el intercambio de
	información con herramientas como DOORS, Jira y Polarion \cite{reqview2026import}.

	\item \textbf{Vistas personalizables:} ofrece tablas configurables con filtros,
	ordenación y columnas personalizadas, permitiendo adaptar la visualización de la
	información a las necesidades de cada proyecto \cite{reqview2026features}.

	\item \textbf{Control de versiones y cambios:} mantiene un historial de
	modificaciones y permite gestionar distintas versiones de los requisitos; al
	trabajar con archivos locales, se integra fácilmente con sistemas como Git y SVN
	\cite{khoroshilov2019}.
\end{itemize}

\subsection{Ventajas}

\begin{itemize}
	\item \textbf{Herramienta ligera y económica:} ReqView es una alternativa más
	sencilla y accesible en comparación con soluciones empresariales como DOORS,
	Polarion o Jama, lo que la hace interesante para pequeñas empresas,
	\textit{startups} y proyectos académicos \cite{reqview2026pricing}.

	\item \textbf{Fácil de aprender y utilizar:} su interfaz es intuitiva y similar a
	una hoja de cálculo, lo que facilita su uso por parte de personas con poca
	experiencia en herramientas de gestión de requisitos \cite{reqview2026version}.

	\item \textbf{Mayor control sobre la información:} al almacenar los datos
	localmente, los usuarios mantienen el control de sus archivos y pueden trabajar
	sin conexión a Internet, sin requerir infraestructura compleja ni servidores
	dedicados \cite{reqview2026features}.
\end{itemize}

\subsection{Limitaciones}

\begin{itemize}
	\item \textbf{Menor cobertura para proyectos críticos:} ReqView no incorpora todas
	las funcionalidades necesarias para proyectos altamente complejos o regulados,
	especialmente en gestión de configuración, análisis de impacto y trabajo
	colaborativo \cite{khoroshilov2019}.

	\item \textbf{Menos adecuado para sectores con alta regulación:} en industrias como
	la aeronáutica, defensa o automoción, donde existen estrictos requisitos de
	certificación, otras herramientas más completas suelen ser una mejor alternativa
	\cite{khoroshilov2019}.

	\item \textbf{Integraciones y funciones avanzadas más limitadas:} aunque permite
	intercambiar información con otras herramientas, no cuenta con el mismo nivel de
	integración ni con las funcionalidades avanzadas de plataformas empresariales como
	IBM DOORS Next o IBM ELM \cite{khoroshilov2019}.
\end{itemize}


% ============================================================
\section{Aplicabilidad al Proyecto Fin de Curso}
% ============================================================

\subsection{Ingeniería de Requisitos y sus Procesos}

La Ingeniería de Requisitos funciona como el proceso que conecta a quienes necesitan un sistema
con quienes lo construyen. Definir requisitos va más allá de registrar lo que el cliente solicita;
el trabajo real consiste en transformar esas necesidades en requisitos concretos, medibles y útiles
para el equipo de desarrollo. El estándar ISO/IEC/IEEE 29148:2018 \cite{iso29148} abarca este
trabajo desde que aparecen las primeras necesidades hasta que los requisitos se validan y se
mantienen vigentes durante toda la vida del sistema.

El SWEBOK \cite{swebok2024} organiza este campo en siete subáreas que incluyen la captura, el
análisis, la especificación y la validación de requisitos, entre otras. Una de las más relevantes
para proyectos en curso es la de consideraciones prácticas, que trata temas como la gestión del
cambio y la trazabilidad. Esto deja en claro que trabajar con requisitos no termina cuando se
levantan: mantenerlos actualizados y rastreables a lo largo del desarrollo tiene tanto peso como
su recopilación inicial. Para MundiPets, los 25 requisitos obtenidos en las entrevistas son el
punto de partida; organizarlos, priorizarlos y vincularlos al desarrollo es la tarea que sigue.

\subsection{Stakeholders como Fuente de Requisitos}

Pohl y Rupp \cite{pohl2015} sostienen que los \textit{stakeholders} representan la fuente más
directa de requisitos en cualquier desarrollo de software. Si algún grupo relevante no se considera
desde el inicio, sus necesidades aparecerán tarde o temprano como cambios durante el proyecto, lo
que implica mayor tiempo y costo. Más allá de recolectar información, el rol del ingeniero de
requisitos incluye gestionar esa relación activamente, entender los objetivos de cada parte y
mediar cuando hay expectativas contradictorias entre sí.

En el caso de MundiPets se trabajó con cinco grupos de \textit{stakeholders}: propietarios de
mascotas, usuarios que buscan adoptar o cruzar, refugios y asociaciones, veterinarias y el equipo
de desarrollo. A cada grupo le correspondió una entrevista estructurada que permitió obtener los
requisitos desde distintos ángulos. Ese proceso de elicitación con múltiples fuentes es
precisamente lo que Pohl y Rupp \cite{pohl2015} describen como una forma sistemática de asegurarse
de que no queden necesidades importantes sin cubrir.

\subsection{Trazabilidad de Requisitos}

Según Pohl y Rupp \cite{pohl2015}, la trazabilidad es la propiedad que permite seguir un requisito
desde donde se originó hasta los artefactos concretos que lo implementan y verifican. Los autores
distinguen entre la trazabilidad hacia el origen del requisito, la trazabilidad hacia los
componentes y casos de prueba que lo realizan, y las relaciones de dependencia entre requisitos.
Cada uno de esos vínculos tiene utilidad práctica distinta dentro del proyecto.

El SWEBOK \cite{swebok2024} señala que sin trazabilidad es difícil analizar el impacto cuando
algo cambia: si un requisito se modifica y no existe registro de qué partes del sistema dependen
de él, el riesgo de introducir errores o dejar funcionalidades incompletas aumenta considerablemente.
En MundiPets esto aplica especialmente a funcionalidades que manejan datos sensibles, como el
historial médico de las mascotas o los procesos de adopción, donde un cambio en un requisito puede
afectar varios puntos del sistema al mismo tiempo.

La norma ISO/IEC/IEEE 29148:2018 \cite{iso29148} trata la trazabilidad como la capacidad de
registrar de dónde viene cada requisito y a qué partes del sistema está asignado, y agrega que
gestionar los requisitos implica procesar los cambios de forma ordenada y mantener a todos los
involucrados al tanto del estado de cada requisito durante el proyecto.

\subsection{Herramientas de Soporte a la Gestión de Requisitos}

Gestionar manualmente la trazabilidad y el seguimiento de requisitos en un equipo de varios
integrantes con decenas de requisitos es una tarea que fácilmente pierde consistencia. El SWEBOK
\cite{swebok2024} señala que disponer de herramientas especializadas facilita actividades como el
control de versiones, el seguimiento de cambios y el registro de artefactos dentro del ciclo de
desarrollo. ZenTao entra en esa categoría: es una plataforma ALM \textit{open source} donde
conviven el \textit{backlog}, los \textit{sprints}, el registro de defectos, las pruebas y la
documentación del proyecto.

Para MundiPets, eso se traduce en algo concreto: los 25 requisitos pasan a ser historias de
usuario, se distribuyen en \textit{sprints}, se vinculan a sus casos de prueba y cualquier error
que aparezca durante el desarrollo queda registrado junto al requisito que lo originó. De esa
forma se construye la trazabilidad post-especificación que Pohl y Rupp \cite{pohl2015} describen:
cada requisito termina enlazado con las tareas que lo desarrollaron y las pruebas que lo verificaron.

Desde los criterios de la norma ISO/IEC/IEEE 29148:2018 \cite{iso29148}, una gestión de requisitos
adecuada requiere que los cambios se controlen y que la información fluya hacia todos los involucrados.
ZenTao responde a eso al guardar el historial de modificaciones por tarea y generar reportes de
avance que sirven como evidencia del progreso, lo cual resulta útil tanto para el trabajo interno
del equipo como para las entregas académicas del curso.

\subsection{Mapeo de Actividades de IR a Módulos de MundiPets}

El SWEBOK V4 organiza la IR en subáreas que incluyen elicitación, análisis, especificación,
validación y gestión de requisitos \cite{swebok2024}. La tabla siguiente muestra cómo ZenTao
da soporte concreto a cada una de esas actividades en el contexto de los módulos de MundiPets,
precisando qué funcionalidad de la herramienta se emplea y a qué módulo del sistema aplica.

\begin{table}[H]
	\caption{Mapeo de actividades de IR con ZenTao para los módulos de MundiPets}
	\label{tab:mapeo-ir}
	\begin{tabularx}{\textwidth}{|l|X|l|}
		\hline
		\textbf{Actividad IR} & \textbf{Uso en ZenTao} & \textbf{Módulo MundiPets} \\
		\hline
		Elicitación & Registro de historias de usuario obtenidas en entrevistas con propietarios de mascotas y refugios & Perfil médico de mascota \\
		\hline
		Especificación & Creación de tareas en el \textit{backlog} con criterios de aceptación y prioridad & Historial de vacunación \\
		\hline
		Trazabilidad & Vinculación de requisitos con casos de prueba y defectos registrados & Proceso de adopción \\
		\hline
		Gestión del cambio & Historial de modificaciones por tarea; notificaciones al equipo ante cambios & Módulo de búsqueda de mascotas \\
		\hline
		Validación & Asociación de casos de prueba directamente con los requisitos correspondientes & Registro de veterinarias \\
		\hline
		Seguimiento de defectos & Registro de errores vinculados al requisito que los originó & Todos los módulos \\
		\hline
	\end{tabularx}
\end{table}

\subsection{Plan de Adopción}

Pohl y Rupp recomiendan adoptar la herramienta de gestión de requisitos antes de que el proyecto
esté en marcha, para evitar los riesgos de introducirla tardíamente \cite{pohl2015}. Para
MundiPets, el plan de adopción de ZenTao contempla las siguientes etapas:

\begin{enumerate}
  \item \textbf{Semana 1:} instalación de ZenTao Open Source en entorno local o en servidor
  compartido del equipo; creación del proyecto ``MundiPets'' y configuración de los módulos
  de \textit{backlog}, \textit{sprints} y pruebas.
  \item \textbf{Semana 2:} migración de los 25 requisitos obtenidos en las entrevistas como
  historias de usuario en el \textit{backlog}, asignando prioridad y responsable a cada una.
  \item \textbf{Semanas 3-4:} vinculación de cada historia de usuario con sus casos de prueba
  correspondientes; establecimiento de los enlaces de trazabilidad entre el módulo de historial
  de vacunación, perfil médico y adopción.
  \item \textbf{Durante el desarrollo:} registro de defectos directamente vinculados al requisito
  que los originó; actualización del historial de cambios ante cualquier modificación de alcance.
\end{enumerate}



\begin{table}[H]
	\caption{Criterios de evaluación de ZenTao para el PFC -- MundiPets}
	\label{tab:criterios}
	\begin{tabularx}{\textwidth}{|l|X|}
		\hline
		\textbf{Criterio} & \textbf{Descripción para el PFC} \\
		\hline
		Costo & Debe ser gratuito o de bajo costo para uso académico \\
		\hline
		Trazabilidad & Debe permitir vincular requisitos con tareas, pruebas y defectos \\
		\hline
		Gestión del cambio & Debe registrar el historial de modificaciones sobre los requisitos \\
		\hline
		Colaboración & Debe soportar trabajo simultáneo de múltiples integrantes del equipo \\
		\hline
		Facilidad de uso & Debe tener una curva de aprendizaje manejable para un equipo académico \\
		\hline
		Integración con el ciclo de desarrollo & Debe soportar \textit{sprints}, \textit{backlog} y seguimiento de defectos \\
		\hline
	\end{tabularx}
\end{table}

\subsection{Comparativa de Aplicabilidad}

\begin{table}[H]
	\caption{Comparativa de funcionalidades de ZenTao frente a herramientas alternativas de IR}
	\label{tab:comparativa}
	\renewcommand{\arraystretch}{1.3}
	\begin{tabularx}{\textwidth}{|l|X|X|X|X|X|}
		\hline
		\textbf{Criterio} & \textbf{ZenTao} & \textbf{Jama Connect} & \textbf{IBM DOORS Next} & \textbf{ReqView} & \textbf{Helix RM} \\
		\hline
		\textbf{Trazabilidad} \cite{hussain2022}
		& Bidireccional: historia $\rightarrow$ tarea $\rightarrow$ caso de prueba $\rightarrow$ defecto
		& Completa con matrices de trazabilidad y vistas de impacto
		& Completa; soporta DO-178C e IEC 62304 para industrias reguladas
		& Básica; trazabilidad a nivel de documento
		& Completa con vistas de árbol de trazabilidad \\
		\hline
		\textbf{Gestión del cambio} \cite{wiegers2013}
		& Historial de modificaciones por tarea; notificaciones automáticas al equipo
		& Control de versiones de requisitos con aprobación formal
		& Control de versiones con flujos de aprobación configurables
		& Historial de revisiones por documento; sin flujos de aprobación
		& Gestión de cambios con flujos de aprobación y auditoría \\
		\hline
		\textbf{Colaboración} \cite{zentao2025}
		& Roles diferenciados (analista, desarrollador, tester); trabajo simultáneo en el mismo proyecto
		& Comentarios inline en requisitos; revisiones colaborativas
		& Roles configurables; integración con Jazz Team Server
		& Edición individual; sin colaboración en tiempo real
		& Revisiones y aprobaciones colaborativas por rol \\
		\hline
		\textbf{Integración con IA} \cite{siddique2023}
		& Sin módulo de IA nativo; integración manual con herramientas externas
		& Asistente de IA para detección de ambigüedades en requisitos
		& Watson AI para análisis de calidad de requisitos (versión enterprise)
		& Sin integración de IA
		& Sin integración de IA nativa documentada \\
		\hline
		\textbf{Costo} \cite{wiegers2013}
		& Gratuito (Open Source); versiones Professional y Enterprise de pago
		& Licencia por usuario; precio desde \$75/mes por usuario
		& Licencia IBM ELM; precio por cotización; alto costo total
		& Gratuito (hasta 40 requisitos); Pro desde \$8/mes
		& Solo por cotización; sin versión gratuita documentada \\
		\hline
		\textbf{Facilidad de uso} \cite{hussain2022}
		& Curva de aprendizaje media; interfaz en inglés y chino; documentación disponible
		& Curva alta; requiere formación para administración
		& Curva muy alta; requiere certificación IBM recomendada
		& Curva baja; interfaz intuitiva y documentación clara
		& Curva alta; orientada a equipos con experiencia en ALM \\
		\hline
		\textbf{Despliegue} \cite{zentao2025}
		& Local (instalación propia) o en la nube
		& Solo en la nube (SaaS)
		& En la nube (IBM ELM) o servidor propio
		& Local (escritorio) o en la nube
		& En la nube o servidor propio \\
		\hline
	\end{tabularx}
\end{table}

El análisis comparativo muestra que ZenTao es la única herramienta del grupo que combina
costo cero, despliegue local sin dependencia de Internet continuo y trazabilidad bidireccional
integrada con gestión de pruebas y defectos en una misma plataforma. Su principal desventaja
frente a Jama Connect e IBM DOORS Next es la ausencia de integración con IA para detección
automática de ambigüedades en requisitos, funcionalidad que Siddique identifica como una
tendencia creciente en las herramientas modernas de IR \cite{siddique2023}. Para el contexto
académico de MundiPets, este déficit no es determinante, dado que el equipo puede realizar
la revisión de calidad de requisitos de forma manual apoyándose en las propiedades definidas
por Pohl y Rupp \cite{pohl2015}.


\newpage

% ============================================================
\section{Vinculación con la Ingeniería de Requisitos}
% ============================================================

La siguiente tabla presenta las conexiones explícitas entre ZenTao y los conceptos centrales
de la Ingeniería de Requisitos estudiados en el curso, siguiendo el framework de Pohl y Rupp,
el SWEBOK V4, la norma ISO/IEC/IEEE 29148:2018 y el modelo de calidad ISO/IEC 25010.

\begin{table}[H]
	\caption{Vinculación de ZenTao con conceptos de la Ingeniería de Requisitos}
	\label{tab:vinculacion-ir}
	\begin{tabularx}{\textwidth}{|l|X|}
		\hline
		\textbf{Concepto de IR} & \textbf{ZenTao se relaciona con este concepto porque\ldots} \\
		\hline
		Trazabilidad (ISO/IEC/IEEE 29148, Pohl y Rupp)
		& permite vincular cada historia de usuario con las tareas de desarrollo,
		los casos de prueba y los defectos que la afectan, satisfaciendo el requisito
		de trazabilidad bidireccional que establece la norma ISO/IEC/IEEE 29148:2018
		\cite{iso29148}. En MundiPets esto se aplicó vinculando HU-02 (historial de
		vacunación) con su caso de prueba CP-02 dentro del módulo de pruebas. \\
		\hline
		Gestión del cambio (SWEBOK V4, subárea Requirements Management)
		& registra el historial de modificaciones de cada requisito y notifica al equipo
		ante cualquier cambio, cubriendo así la subárea de gestión de requisitos del
		SWEBOK V4, que señala que controlar los cambios y mantener la trazabilidad son
		actividades tan importantes como la elicitación inicial \cite{swebok2024}. \\
		\hline
		Propiedades de los requisitos: verificabilidad y trazabilidad
		& al asociar cada historia de usuario con un caso de prueba concreto, ZenTao
		operacionaliza dos propiedades que Pohl y Rupp consideran fundamentales en
		cualquier requisito bien formado: la verificabilidad, que exige que exista al
		menos un método para comprobar si el requisito se cumple, y la trazabilidad,
		que exige que cada requisito pueda rastrearse hasta su origen y hasta los
		artefactos que lo implementan \cite{pohl2015}. \\
		\hline
		Actividades core de IR: especificación y validación (SWEBOK V4)
		& el módulo \textit{backlog} de ZenTao soporta la especificación de requisitos
		en formato de historias de usuario con criterios de aceptación en sintaxis
		\textit{Given-When-Then}; el módulo de pruebas soporta la validación al
		permitir ejecutar y registrar los resultados de cada caso de prueba vinculado
		a un requisito \cite{swebok2024}. Ambas son actividades core de la IR según
		el SWEBOK V4. \\
		\hline
		Atributos de calidad del software (ISO/IEC 25010)
		& ZenTao contribuye directamente a dos atributos de calidad del modelo
		ISO/IEC 25010: la \textbf{mantenibilidad}, al mantener un historial de cambios
		que facilita la comprensión y modificación del sistema, y la
		\textbf{trazabilidad funcional}, al garantizar que cada requisito pueda
		rastrearse hasta su verificación \cite{iso25010}. \\
		\hline
		Colaboración entre stakeholders (Pohl y Rupp, framework de IR)
		& centraliza la información del proyecto en un repositorio único accesible
		para todos los roles del equipo (analista, desarrollador, tester), lo que
		responde a la necesidad de gestionar activamente la relación con los
		\textit{stakeholders} que Pohl y Rupp identifican como una responsabilidad
		central del ingeniero de requisitos \cite{pohl2015}. En MundiPets esto
		permite que los tres integrantes del equipo trabajen sobre los mismos
		requisitos sin inconsistencias. \\
		\hline
	\end{tabularx}
\end{table}

% ============================================================
\section{Análisis Crítico}
% ============================================================

\subsection{¿ZenTao reemplaza o complementa al analista de requisitos?}

ZenTao complementa al analista de requisitos; no lo reemplaza. La herramienta automatiza
actividades de registro, seguimiento y notificación de cambios, pero las tareas de mayor
valor cognitivo de la IR, como la elicitación de necesidades implícitas, la negociación entre
\textit{stakeholders} con intereses contradictorios y la validación del significado real de un
requisito, continúan dependiendo del criterio del analista. Hussain y Mkpojiogu señalan que
las herramientas ALM de propósito general gestionan la información de los requisitos, pero no
pueden sustituir el razonamiento que se necesita para transformar una necesidad vaga en un
requisito verificable \cite{hussain2022}. En MundiPets, ZenTao organiza los 25 requisitos
obtenidos en entrevistas, pero la decisión de cómo priorizar, qué alcance tiene cada uno y
cómo resolver conflictos entre propietarios de mascotas y refugios es responsabilidad del equipo,
no de la herramienta.

\subsection{¿Qué propiedad de calidad mejora y cuál compromete?}

ZenTao mejora principalmente la mantenibilidad y la trazabilidad
del sistema de requisitos: el historial de cambios por tarea y los vínculos entre historias de
usuario, casos de prueba y defectos facilitan comprender qué partes del sistema se ven afectadas
ante cualquier modificación, atributos directamente contemplados en el modelo ISO/IEC 25010
\cite{iso25010}.

Sin embargo, ZenTao puede comprometer la portabilidad y la independencia
de datos: al almacenar toda la información del proyecto en su base de datos MySQL propia,
la exportación de requisitos hacia otras herramientas no es directa ni está estandarizada. Si el
equipo decide migrar a otra plataforma, los vínculos de trazabilidad construidos en ZenTao
no se transfieren automáticamente en formatos estándar como ReqIF, a diferencia de herramientas
como ReqView o IBM DOORS Next que soportan este estándar de intercambio de forma nativa.
Wiegers y Beatty advierten que la dependencia de un formato propietario es un riesgo que
debe considerarse al evaluar una herramienta de gestión de requisitos \cite{wiegers2013}.

\subsection{¿Es aplicable en el contexto local de la UTEQ?}

Sí, ZenTao es aplicable en el contexto académico de la UTEQ por tres razones concretas. Primero,
su edición \textit{Open Source} no requiere licencia comercial, lo que elimina la principal barrera
económica para equipos de estudiantes. Segundo, su instalación es local y no depende de
conectividad continua a Internet, lo que la hace viable en entornos con acceso irregular a la red.
Tercero, la interfaz en inglés, aunque puede representar una dificultad inicial, es manejable para
estudiantes de ingeniería de software en el nivel de cuarto semestre. Pohl y Rupp señalan que
al evaluar la viabilidad de una herramienta en un contexto específico, el costo de adopción total,
que incluye tiempo de aprendizaje, infraestructura y soporte, es tan determinante como sus
funcionalidades técnicas \cite{pohl2015}. En ese balance, ZenTao es la opción más favorable
frente a Jama Connect, IBM DOORS Next o Helix RM para el contexto de la carrera.

\subsection{¿Qué riesgo ético identifica?}

El principal riesgo ético en el uso de ZenTao para MundiPets es la privacidad de los datos
médicos de las mascotas. El sistema almacena información sensible: historial de vacunación,
enfermedades crónicas, estado reproductivo y antecedentes genéticos. Si ZenTao se despliega
en un servidor compartido sin controles de acceso adecuados, cualquier integrante del equipo o
persona con acceso al servidor podría consultar o modificar esos datos sin autorización. La norma
ISO/IEC/IEEE 29148:2018 establece que los requisitos de privacidad y seguridad deben ser
identificados y gestionados desde las etapas tempranas del desarrollo \cite{iso29148}, lo que
implica que el equipo debe definir explícitamente qué roles tienen permiso de lectura y escritura
sobre cada módulo de MundiPets dentro de ZenTao, antes de registrar información real de mascotas
o usuarios en la plataforma. Un segundo riesgo es la dependencia tecnológica: si el
proyecto académico migra a producción real, ZenTao Open Source no ofrece garantías de soporte
a largo plazo, lo que podría dejar al sistema sin herramienta de gestión de requisitos en etapas
críticas del desarrollo.

\newpage

\section{Conclusiones}

ZenTao constituye una herramienta útil para apoyar actividades relacionadas con la Ingeniería de
Requisitos, especialmente aquellas vinculadas con la trazabilidad, la gestión del cambio y la
colaboración entre \textit{stakeholders} \cite{zentao2025}. Su modelo \textit{open source} y su
integración de múltiples módulos dentro de una misma plataforma la convierten en una alternativa
adecuada para proyectos académicos y organizaciones de pequeña y mediana escala.

Sin embargo, sus limitaciones evidencian que ninguna herramienta puede reemplazar completamente
el trabajo del analista de requisitos \cite{hussain2022}. ZenTao debe considerarse un instrumento
de apoyo y no un sustituto del criterio humano, ya que las tareas de elicitación, negociación y
validación continúan dependiendo de la capacidad del equipo para comprender las necesidades de los
\textit{stakeholders} y mediar entre expectativas distintas \cite{pohl2015}.

En el contexto del PFC MundiPets, ZenTao representa la alternativa más viable entre las herramientas
estudiadas, dado que satisface los criterios de costo, trazabilidad, gestión del cambio y
colaboración sin requerir licencias comerciales. La calidad de los requisitos, sin embargo,
continúa dependiendo de los procesos de elicitación y validación que el equipo lleva a cabo fuera
de la plataforma, en concordancia con los principios establecidos en el SWEBOK V4 \cite{swebok2024}
y en la norma ISO/IEC/IEEE 29148 \cite{iso29148}.

Desde la perspectiva de la calidad del software, ZenTao contribuye principalmente a mejorar la
mantenibilidad y la trazabilidad de los requisitos, atributos alineados con el modelo de calidad
ISO/IEC 25010 \cite{iso25010}.

% ============================================================
\newpage
\printbibliography[heading=bibnumbered, title={Referencias}]

\newpage

% ============================================================

\section{Anexos}

\subsection{Manual de Usuario: ZenTao Open Source}
\label{app:manual}

Esta sección documenta los pasos para instalar ZenTao Open Source y configurar un proyecto
con requisitos del sistema MundiPets, siguiendo la metodología de gestión de requisitos
descrita en el SWEBOK V4 \cite{swebok2024}.

\subsubsection{Instalación}

ZenTao Open Source puede instalarse mediante el paquete de instalación multiplataforma
disponible en \url{https://www.zentao.net/download/pms22.2-86288.html}. Los pasos generales son:

\begin{enumerate}
	\item Descargar el instalador para el sistema operativo correspondiente (Windows, Linux o macOS).
	\item Ejecutar el instalador; ZenTao despliega automáticamente un servidor
	Apache, PHP y MySQL sin necesidad de configuración adicional.
	\item Acceder desde el navegador a \texttt{http://localhost/zentao} y completar la configuración
	inicial: nombre de la empresa, usuario administrador y contraseña.
	\item Crear el proyecto ``MundiPets'' desde el menú \textbf{Proyectos > Crear proyecto},
	seleccionando metodología \textit{Scrum} y definiendo la duración del semestre como horizonte
	temporal.
\end{enumerate}

% IMAGEN 1: zentao_instalador.png
% Es la pantalla azul del asistente de instalación con el logo de ZenTao y el botón 安装
\begin{figure}[H]
	\centering
	\fbox{\includegraphics[width=0.6\textwidth]{zentao_instalador.png}}
	\caption{Asistente de instalación de ZenTao Open Source: selección de directorio de
		instalación con el servidor Apache, PHP y MySQL integrados.}
	\label{fig:zentao-instalador}
\end{figure}

% IMAGEN 2: zentao_panel_servicios.png
% Es el panel de control con los servicios Apache, MySQL, etc. y el botón "Visit ZenTao"
\begin{figure}[H]
	\centering
	\fbox{\includegraphics[width=\textwidth]{zentao_panel_servicios.png}}
	\caption{Panel de servicios de ZenTao Integrated Environment: estado de los servicios
		Apache, MySQL, PHP y Redis tras la instalación exitosa. Desde aquí se accede
		a la interfaz web mediante el botón \textit{Visit ZenTao}.}
	\label{fig:zentao-panel-servicios}
\end{figure}

% IMAGEN 3: zentao_bienvenida.png
% Es la pantalla de bienvenida en localhost con el logo de ZenTao y las credenciales iniciales
\begin{figure}[H]
	\centering
	\fbox{\includegraphics[width=\textwidth]{zentao_bienvenida.png}}
	\caption{Pantalla de bienvenida de ZenTao en \texttt{http://localhost}: confirmación de
		instalación exitosa con las credenciales iniciales del administrador.}
	\label{fig:zentao-bienvenida}
\end{figure}

\newpage

\subsubsection{Panel principal y proyecto MundiPets}

Una vez completada la instalación e iniciada sesión con las credenciales de administrador,
ZenTao presenta un panel central que centraliza toda la actividad del proyecto. Desde este
panel es posible acceder a los módulos de \textit{Program}, \textit{Product},
\textit{Project}, \textit{Execution} y \textit{Test}, que corresponden a las etapas del
ciclo de vida del desarrollo descritas en el SWEBOK V4 \cite{swebok2024}. En la sección
\textit{My Pending Reviews} se visualizan los requisitos pendientes de revisión del producto
\textit{SistemaMudiPets}, lo que permite al equipo mantener un control continuo sobre el
estado de cada historia de usuario registrada.

% IMAGEN 4: interfaz_motor.png
% Es el Dashboard de ZenTao ya logueado, se ve SistemamudiPets en la columna Product
\begin{figure}[H]
	\centering
	\fbox{\includegraphics[width=\textwidth]{interfaz_motor.png}}
	\caption{Panel principal (Dashboard) de ZenTao con el proyecto MundiPets activo:
		se visualizan los módulos de revisiones pendientes, historias, bugs y tareas
		asignadas al equipo dentro del producto \textit{SistemaMudiPets}.}
	\label{fig:zentao-dashboard}
\end{figure}

\newpage

\subsubsection{Registro de Historias de Usuario (Requisitos MundiPets)}

Una vez creado el proyecto, los requisitos obtenidos en las entrevistas se registran como
historias de usuario en el módulo \textbf{Product \textgreater{} Stories}:

\begin{enumerate}
	\item Acceder a \textbf{Product \textgreater{} Stories \textgreater{} Create Story}.
	\item Ingresar el título de la historia, por ejemplo: \textit{``Como propietario de mascota,
		quiero registrar el historial de vacunación de mi mascota para consultarlo en cualquier
		momento.''} (requisito funcional HU-02, módulo Historial de Vacunación).
	\item Asignar prioridad (alta/media/baja) y el responsable del equipo.
	\item Añadir los criterios de aceptación en el campo de descripción, en formato
	\textit{Given-When-Then}.
	\item Repetir el proceso para los 25 requisitos obtenidos de los cinco grupos de
	\textit{stakeholders} de MundiPets.
\end{enumerate}

% ---------------------------------------------------------------
% IMAGEN 5 (PENDIENTE): zentao_interfaz.png
% Necesitas tomar esta captura:
%   Ve a Product > SistemamudiPets > Stories
%   Deben verse varias historias de usuario del PFC en la lista.
%   Asegúrate de que la columna "Title" muestre nombres relacionados
%   con MundiPets (historial de vacunación, adopción, veterinarias, etc.)
%   y que se vea la columna de prioridad y responsable.
% ---------------------------------------------------------------
\begin{figure}[H]
	\centering
	\fbox{\includegraphics[width=\textwidth]{zentao_interfaz.png}}
	\caption{Módulo de historias de usuario del producto \textit{SistemaMudiPets} en ZenTao:
		listado de requisitos funcionales con prioridad y responsable asignado, en concordancia
		con la gestión de requisitos descrita en la norma ISO/IEC/IEEE 29148:2018 \cite{iso29148}.}
	\label{fig:zentao-stories}
\end{figure}

\newpage

\subsubsection{Vinculación de Requisitos con Casos de Prueba}

Para establecer la trazabilidad entre requisitos y pruebas, ZenTao permite vincular cada
historia de usuario con los casos de prueba que la verifican:

\begin{enumerate}
	\item Desde \textbf{Test \textgreater{} Cases}, crear un caso de prueba asociado a una historia
	de usuario del producto \textit{SistemaMudiPets}.
	\item En el campo \textbf{Related Story}, seleccionar el requisito correspondiente
	(por ejemplo HU-13, módulo Proceso de Adopción).
	\item ZenTao registra el vínculo; desde la vista de la historia es posible verificar
	qué requisitos tienen cobertura de prueba y cuáles no.
\end{enumerate}

% ---------------------------------------------------------------
% IMAGEN 6 (PENDIENTE): zentao_trazabilidad.png
% Necesitas tomar esta captura:
%   Ve a Test > Cases y abre el detalle de un caso de prueba.
%   Debe verse el campo "Related Story" con una historia de MundiPets vinculada.
%   O bien, abre una historia de usuario y muestra la pestaña de casos de prueba
%   relacionados en la parte inferior.
% ---------------------------------------------------------------
\begin{figure}[H]
	\centering
	\fbox{\includegraphics[width=\textwidth]{zentao_trazabilidad.png}}
	\caption{Vista de trazabilidad en ZenTao: caso de prueba CP-13 del módulo Proceso de Adopción
		vinculado a la historia de usuario HU-13 en \textit{SistemaMudiPets},
		permitiendo verificar la cobertura de cada requisito del PFC \cite{swebok2024}.}
	\label{fig:zentao-trazabilidad}
\end{figure}

\newpage

\subsection{GitHub}

Enlace al repositorio de GitHub con los aportes individuales de los integrantes acerca de la investigación: \url{https://github.com/MiloSaurio4kHD/ExposicionZenTao/tree/main}

\begin{figure}[H]
	\centering
	\fbox{\includegraphics[width=\textwidth]{repositorio.png}}
	\caption{Repositorio del informe de ZenTao}
	\label{fig:repositorio-github}
\end{figure}

\begin{figure}[H]
	\centering
	\fbox{\includegraphics[width=\textwidth]{commits.png}}
	\caption{Commits al repositorio del informe de ZenTao}
	\label{fig:commit-github}
\end{figure}


\end{document}
